\Article{240}
\textsc{Satz.} \textit{Aus den Formen $f,f'$ sei die Form $F$ zusammengesetzt; aus $F$ und $f''$ die Form $\mathfrak{F}$; aus $f,f''$ die Form $F'$; aus $F'$ und $f'$ die Form $\mathfrak{F}'$. Dann sind die Formen $\mathfrak{F},\mathfrak{F}'$ eigentlich äquivalent.}

\textit{Beweis.} I. Es sei
\[
\begin{aligned}
f&=axx+2bxy+cyy\\
f'&=a'x'x'+2b'x'y'+c'y'y'\\
f''&=a''x''x''+2b''x''y''+c''y''y''\\
F&=AXX+2BXY+CYY\\
F'&=A'X'X'+2B'X'Y'+C'Y'Y'\\
\mathfrak{F}&=\mathfrak{A}\mathfrak{X}\mathfrak{X}+2\mathfrak{B}\mathfrak{X}\mathfrak{Y}+\mathfrak{C}\mathfrak{Y}\mathfrak{Y}\\
\mathfrak{F}'&=\mathfrak{A}'\mathfrak{X}'\mathfrak{X}'+2\mathfrak{B}'\mathfrak{X}'\mathfrak{Y}'+\mathfrak{C}'\mathfrak{Y}'\mathfrak{Y}'
\end{aligned}
\]
Die Determinanten dieser sieben Formen seien jeweils $d,d',d'',D,D',\mathfrak{D},\mathfrak{D}'$, die alle dasselbe Vorzeichen haben und im Verhältnis von Quadraten zueinander stehen. Weiter sei $m$ der größte gemeinsame Teiler der Zahlen $a,2b,c$, und $m',m'',M$ sollen für die Formen $f',f'';F$ eine entsprechende Bedeutung haben. Dann ist nach dem vierten Schluss des Artikels 235 $D$ der größte gemeinsame Teiler der Zahlen $dm'm',d'mm$, also $Dm''m''$ der größte gemeinsame Teiler der Zahlen $dm'm'm''m'',d'mmm''m''$; $M=mm'$; $\mathfrak{D}$ ist der größte gemeinsame Teiler der Zahlen $Dm''m''$, $d''MM$, also der Zahlen $Dm''m'',d''mmm'm'$. Daraus folgt, dass $\mathfrak{D}$ der größte gemeinsame Teiler der drei Zahlen $dm'm'm''m'',d'mmm''m'',d''mmm'm'$ ist; ebenso ist aber $\mathfrak{D}'$ der größte gemeinsame Teiler derselben drei Zahlen. Da $\mathfrak{D},\mathfrak{D}'$ dasselbe Vorzeichen haben, ist also $\mathfrak{D}=\mathfrak{D}'$, oder die Formen $\mathfrak{F},\mathfrak{F}'$ haben dieselbe Determinante.

\SourceLeaf{261}
II. Nun gehe $F$ in $ff'$ durch die Substitution
\[
\begin{aligned}
X&=pxx'+p'xy'+p''yx'+p'''yy'\\
Y&=qxx'+q'xy'+q''yx'+q'''yy'
\end{aligned}
\]
und $\mathfrak{F}$ in $Ff''$ durch die Substitution
\[
\begin{aligned}
\mathfrak{X}&=\mathfrak{p}Xx''+\mathfrak{p}'Xy''+\mathfrak{p}''Yx''+\mathfrak{p}'''Yy''\\
\mathfrak{Y}&=\mathfrak{q}Xx''+\mathfrak{q}'Xy''+\mathfrak{q}''Yx''+\mathfrak{q}'''Yy''
\end{aligned}
\]
über, und wir bezeichnen die positiven Quadratwurzeln aus $\frac{d}{D},\frac{d'}{D},\frac{D}{\mathfrak{D}}.\ \frac{d''}{\mathfrak{D}}$ mit $n,n',\mathfrak{N},\mathfrak{n}''$. Dann erhalten wir nach Artikel 235 achtzehn Gleichungen, von denen die eine Hälfte zur Transformation der Form $F$ in $ff'$ gehört, die andere zur Transformation der Form $\mathfrak{F}$ in $Ff''$. Die erste ist $pq'-qp'=an'$; nach ihrem Muster lassen sich die übrigen leicht bilden, die wir hier der Kürze wegen weglassen. Die Größen $n,n',\mathfrak{N},\mathfrak{n}''$ sind zwar rational, aber nicht notwendig ganze Zahlen.

III. Setzt man die Werte von $X,Y$ in die Werte von $\mathfrak{X},\mathfrak{Y}$ ein, so ergibt sich eine Substitution der folgenden Art:
\[
\begin{aligned}
\mathfrak{X}={}&(1)xx'x''+(2)xx'y''+(3)xy'x''+(4)xy'y''\\
&+(5)yx'x''+(6)yx'y''+(7)yy'x''+(8)yy'y''\\
\mathfrak{Y}={}&(9)xx'x''+(10)xx'y''+(11)xy'x''+(12)xy'y''\\
&+(13)yx'x''+(14)yx'y''+(15)yy'x''+(16)yy'y''
\end{aligned}
\]
durch die $\mathfrak{F}$ offenbar in das Produkt $ff'f''$ übergeht. Der Koeffizient $(1)$ ist $=p\mathfrak{p}+q\mathfrak{p}''$; die Werte der übrigen fünfzehn geben wir nicht an, da sie jeder ohne Schwierigkeit entwickeln kann. Bezeichnen wir die Zahl $(1)(10)-(2)(9)$ mit $(1,2)$, die Zahl $(1)(11)-(3)(9)$ mit $(1,3)$ und allgemein $(g)(8+h)-(h)(8+g)$ mit $(g,h)$, wobei $g,h$ verschiedene ganze Zahlen zwischen 1 und 16 sind und die größere $h$ ist\HistoricalNote{*)}{Die jetzige Bedeutung dieser Zeichen darf nicht mit der verwechselt werden, die sie in Artikel 234 hatten; denn die durch diese Zeichen \textit{hier} ausgedrückten Zahlen entsprechen genau denen, die in Artikel 234 mit den Zahlen multipliziert wurden, die \textit{dort} durch ähnliche Zeichen bezeichnet waren.}; auf diese Weise erhalten wir insgesamt achtundzwanzig Zeichen. Bezeichnen wir nun die positiven Quadratwurzeln aus $\frac{d}{\mathfrak{D}},\frac{d'}{\mathfrak{D}}$ mit $\mathfrak{n},\mathfrak{n}'$ (die $=n\mathfrak{N},n'\mathfrak{N}$ sind), so ergeben sich die folgenden 28 Gleichungen:

\SourceLeaf{262}
\[
\begin{aligned}
(1,2)&=aa'\mathfrak{n}''\\
(1,3)&=aa''\mathfrak{n}'\\
(1,4)&=ab'\mathfrak{n}''+ab''\mathfrak{n}'\\
(1,5)&=a'a''\mathfrak{n}\\
(1,6)&=a'b\mathfrak{n}''+a'b''\mathfrak{n}\\
(1,7)&=a''b\mathfrak{n}'+a''b'\mathfrak{n}\\
(1,8)&=bb'\mathfrak{n}''+bb''\mathfrak{n}'+b'b''\mathfrak{n}+\mathfrak{D}\mathfrak{n}\mathfrak{n}'\mathfrak{n}'',\\
(2,3)&=ab''\mathfrak{n}'-ab'\mathfrak{n}''\\
(2,4)&=ac''\mathfrak{n}'\\
(2,5)&=a'b''\mathfrak{n}-a'b\mathfrak{n}''\\
(2,6)&=a'c''\mathfrak{n}\\
(2,7)&=bb''\mathfrak{n}'+b'b''\mathfrak{n}-bb'\mathfrak{n}''-\mathfrak{D}\mathfrak{n}\mathfrak{n}'\mathfrak{n}'',\\
(2,8)&=bc''\mathfrak{n}'+b'c''\mathfrak{n}\\
(3,4)&=ac'\mathfrak{n}''
\end{aligned}
\]
\[
\begin{aligned}
(3,5)&=a''b'\mathfrak{n}-a''b\mathfrak{n}'\\
(3,6)&=bb'\mathfrak{n}''+b'b''\mathfrak{n}-bb''\mathfrak{n}'-\mathfrak{D}\mathfrak{n}\mathfrak{n}'\mathfrak{n}''\\
(3,7)&=a''c'\mathfrak{n}\\
(3,8)&=bc'\mathfrak{n}''+b''c\mathfrak{n}\\
(4,5)&=b'b''\mathfrak{n}-bb'\mathfrak{n}''-bb''\mathfrak{n}'+\mathfrak{D}\mathfrak{n}\mathfrak{n}'\mathfrak{n}''\\
(4,6)&=b'c''\mathfrak{n}-bc''\mathfrak{n}'\\
(4,7)&=b''c'\mathfrak{n}-bc'\mathfrak{n}''\\
(4,8)&=c'c''\mathfrak{n}\\
(5,6)&=ca'\mathfrak{n}''\\
(5,7)&=ca''\mathfrak{n}'\\
(5,8)&=b'c\mathfrak{n}''+b''c\mathfrak{n}'\\
(6,7)&=b''c\mathfrak{n}'-b'c\mathfrak{n}''\\
(6,8)&=cc''\mathfrak{n}'\\
(7,8)&=cc'\mathfrak{n}''
\end{aligned}
\]
die wir mit $\Phi$ bezeichnen werden, und neun weitere:
\[
\begin{aligned}
(10)(11)-(9)(12)&=a\mathfrak{n}'\mathfrak{n}''\mathfrak{A}\\
(1)(12)-(2)(11)-(3)(10)+(4)(9)&=2a\mathfrak{n}'\mathfrak{n}''\mathfrak{B}\\
(2)(3)-(1)(4)&=a\mathfrak{n}'\mathfrak{n}''\mathfrak{C}\\
-(9)(16)+(10)(15)+(11)(14)-(12)(13)&=2b\mathfrak{n}'\mathfrak{n}''\mathfrak{A}\\
\begin{gathered}
(1)(16)-(2)(15)-(3)(14)+(4)(13)\{\\
+(5)(12)-(6)(11)-(7)(10)+(8)(9)\}
\end{gathered}&=4b\mathfrak{n}'\mathfrak{n}''\mathfrak{B}\\
-(1)(8)+(2)(7)+(3)(6)-(4)(5)&=2b\mathfrak{n}'\mathfrak{n}''\mathfrak{C}\\
(14)(15)-(13)(16)&=c\mathfrak{n}'\mathfrak{n}''\mathfrak{A}\\
(5)(16)-(6)(15)-(7)(14)+(8)(13)&=2c\mathfrak{n}'\mathfrak{n}''\mathfrak{B}\\
(6)(7)-(5)(8)&=c\mathfrak{n}'\mathfrak{n}''\mathfrak{C}
\end{aligned}
\]
die wir mit $\Psi$ bezeichnen\HistoricalNote{*)}{Man beachte, dass sich 18 weitere, diesen $\Psi$ ähnliche Gleichungen herleiten lassen, bei denen rechts anstelle der Faktoren $a$, $2b$, $c$ die Faktoren $a'$, $2b'$, $c'$; $a''$, $2b''$, $c''$ stehen. Da wir sie für unseren Zweck nicht brauchen, lassen wir sie weg.}.

IV. Die Entstehung aller dieser 37 Gleichungen herzuleiten, wäre zu lang. Es genügt, einige zu bestätigen, nach deren Muster sich die übrigen ohne große Schwierigkeit beweisen lassen.

\leavevmode\SourceLeaf{263}
1) Es ist
\[
\begin{aligned}
(1,2)&=(1)(10)-(2)(9)\\
&=(\mathfrak{p}\mathfrak{q}'-\mathfrak{q}\mathfrak{p}')pp\\&\quad+(\mathfrak{p}\mathfrak{q}'''-\mathfrak{q}\mathfrak{p}'''-\mathfrak{p}'\mathfrak{q}''+\mathfrak{q}'\mathfrak{p}'')pq\\&\quad+(\mathfrak{p}''\mathfrak{q}'''-\mathfrak{q}''\mathfrak{p}''')qq\\
&=\mathfrak{n}''(App+2Bpq+Cqq)=\mathfrak{n}''aa'
\end{aligned}
\]
die erste Gleichung.

2) Es ergibt sich
\[
\begin{aligned}(1,3)&=(1)(11)-(3)(9)\\&=(\mathfrak{p}\mathfrak{q}''-\mathfrak{q}\mathfrak{p}'')(pq'-qp')\\&=a''\mathfrak{N}an'=aa''\mathfrak{n}'\end{aligned}
\]
die zweite Gleichung.

3) Es ist
\[
\begin{aligned}
(1,8)&=(1)(16)-(8)(9)\\
&=(\mathfrak{p}\mathfrak{q}'-\mathfrak{q}\mathfrak{p}')pp'''\\&\quad+(\mathfrak{p}\mathfrak{q}'''-\mathfrak{q}\mathfrak{p}''')pq'''\\&\quad-(\mathfrak{p}'\mathfrak{q}''-\mathfrak{q}'\mathfrak{p}'')qp'''\\&\quad+(\mathfrak{p}''\mathfrak{q}'''-\mathfrak{q}''\mathfrak{p}''')qq'''\\
&=\mathfrak{n}''(App'''+B(pq'''+qp''')+Cqq''')\\&\quad+b''\mathfrak{N}(pq'''-qp''')\\
&=\mathfrak{n}''(bb'+\surd dd')+b''\mathfrak{N}(b'n+bn')\text{\textsuperscript{*)}}\\
&=\mathfrak{n}''bb'+\mathfrak{n}'bb''+\mathfrak{n}b'b''+\mathfrak{D}\mathfrak{n}\mathfrak{n}'\mathfrak{n}''
\end{aligned}
\]
\HistoricalDisplayNote{*)}{Das folgt aus Gleichung 10 des Artikels 235 und dem Folgenden. Die Wurzelgröße $\surd dd'$ wird $=Dnn'=\mathfrak{D}nn'\mathfrak{N}\mathfrak{N}=\mathfrak{D}\mathfrak{n}\mathfrak{n}'$.}
die achte Gleichung in $\Phi$. Die übrigen Gleichungen überlassen wir den Lesern zur Bestätigung.

V. Aus den Gleichungen $\Phi$ folgt auf die folgende Weise, dass die achtundzwanzig Zahlen $(1,2)$, $(1,3)$ usw. keinen gemeinsamen Teiler haben. Zuerst betrachten wir die siebenundzwanzig Produkte aus je drei Faktoren, deren erster entweder $\mathfrak{n}$, deren zweiter eine der Zahlen $a'$, $2b'$, $c'$ und deren dritter eine der Zahlen $a''$, $2b''$, $c''$ ist; oder deren erster $\mathfrak{n}'$, deren zweiter eine der Zahlen $a$, $2b$, $c$ und deren dritter eine der Zahlen $a''$, $2b''$, $c''$ ist; oder schließlich deren erster $\mathfrak{n}''$, deren zweiter eine der Zahlen $a$, $2b$, $c$ und deren dritter eine der Zahlen $a'$, $2b'$, $c'$ ist---jedes dieser siebenundzwanzig Produkte ist nach den Gleichungen $\Phi$ entweder einer der achtundzwanzig Zahlen $(1,2)$, $(1,3)$ usw. oder der Summe oder Differenz mehrerer davon gleich (\textit{zum Beispiel} $\mathfrak{n}a'a''=(1,5)$, $2\mathfrak{n}a'b''=(1,6)+(2,5)$, $4\mathfrak{n}b'b''=(1,8)+(2,7)+(3,6)+(4,5)$ und ebenso bei den übrigen). Hätten diese Zahlen also einen gemeinsamen Teiler, so müsste dieser auch alle jene Produkte teilen. Daraus folgt mithilfe des Artikels 40 und nach der schon öfter benutzten Methode leicht, dass derselbe Teiler auch die Zahlen $\mathfrak{n}m'm''$, $\mathfrak{n}'mm''$, $\mathfrak{n}''mm'$ teilen müsste, also deren
\SourceLeaf{264}
Quadrate, nämlich $\frac{dm'm'm''m''}{\mathfrak{D}}$, $\frac{d'mmm''m''}{\mathfrak{D}}$, $\frac{d''mmm'm'}{\mathfrak{D}}$, durch sein Quadrat teilbar wären, \textit{was widersinnig ist}, weil nach I der größte gemeinsame Teiler der drei Zähler $\mathfrak{D}$ ist und somit die Quadrate selbst keinen gemeinsamen Teiler haben können.

VI. Das alles betrifft die Transformation der Form $\mathfrak{F}$ in $ff'f''$ und wurde aus den Transformationen der Form $F$ in $ff'$ und der Form $\mathfrak{F}$ in $Ff''$ hergeleitet. Ganz ebenso ergibt sich aber aus den Transformationen der Form $F'$ in $ff''$ und der Form $\mathfrak{F}'$ in $F'f'$ eine Transformation der Form $\mathfrak{F}'$ in $ff'f''$ der folgenden Art:
\[
\begin{aligned}
\mathfrak{X}'&=(1)'xx'x''+(2)'xx'y''+(3)'xy'x''+\text{etc.}\\
\mathfrak{Y}'&=(9)'xx'x''+(10)'xx'y''+(11)'xy'x''+\text{etc.}
\end{aligned}
\]
(wobei wir alle Koeffizienten wie bei der Transformation der Form $\mathfrak{F}$ in $ff'f''$ bezeichnen und zur Unterscheidung jedem einen Strich anfügen). Daraus folgen wie zuvor 28 den Gleichungen $\Phi$ entsprechende Gleichungen, die wir mit $\Phi'$ bezeichnen, und neun weitere, den Gleichungen $\Psi$ entsprechende, die wir mit $\Psi'$ bezeichnen. Bezeichnen wir nämlich
\[
(1)'(10)'-(2)'(9)'\text{ mit }(1,2)',\quad
(1)'(11)'-(3)'(9)'\text{ mit }(1,3)'
\]
so sind die Gleichungen $\Phi'$
\[
(1,2)'=aa'\mathfrak{n}'',\quad (1,3)'=aa''\mathfrak{n}',\quad\text{etc.}
\]
und die Gleichungen $\Psi'$
\[
(10)'(11)'-(9)'(12)'=a\mathfrak{n}'\mathfrak{n}''\mathfrak{A}',\quad\text{etc.}
\]
(Die ausführlichere Entwicklung überlassen wir der Kürze wegen den Lesern; Kenner werden übrigens finden, dass eine neue Rechnung nicht einmal nötig ist, sondern sich die erste Analyse durch Analogie leicht hierher übertragen lässt.) Danach folgt aus $\Phi$ und $\Phi'$ sofort
\[
(1,2)=(1,2)',\quad (1,3)=(1,3)',\quad (1,4)=(1,4)',\quad (2,3)=(2,3)',\quad\text{etc.}
\]
Daraus und weil alle $(1,2)$, $(1,3)$, $(2,3)$ usw. nach V keinen gemeinsamen Teiler haben, folgt mithilfe des Lemmas aus Artikel 234, dass sich vier ganze Zahlen $\alpha$, $\beta$, $\gamma$, $\delta$ so bestimmen lassen, dass
\[
\begin{aligned}
\alpha(1)'+\beta(9)'&=(1),\\\alpha(2)'+\beta(10)'&=(2),\\\alpha(3)'+\beta(11)'&=(3),\quad\text{etc.}\\
\gamma(1)'+\delta(9)'&=(9),\\\gamma(2)'+\delta(10)'&=(10),\\\gamma(3)'+\delta(11)'&=(11),\quad\text{etc.}
\end{aligned}
\]
und $\alpha\delta-\beta\gamma=1$ gilt.

\SourceLeaf{265}
VII. Daraus bestätigt man leicht, indem man aus den ersten drei Gleichungen $\Psi$ die Werte von $a\mathfrak{A}$, $a\mathfrak{B}$, $a\mathfrak{C}$ und aus den ersten drei Gleichungen $\Psi'$ die Werte von $a\mathfrak{A}'$, $a\mathfrak{B}'$, $a\mathfrak{C}'$ einsetzt, dass
\[
\begin{aligned}
a(\mathfrak{A}\alpha\alpha+2\mathfrak{B}\alpha\gamma+\mathfrak{C}\gamma\gamma)&=a\mathfrak{A}'\\
a(\mathfrak{A}\alpha\beta+\mathfrak{B}(\alpha\delta+\beta\gamma)+\mathfrak{C}\gamma\delta)&=a\mathfrak{B}'\\
a(\mathfrak{A}\beta\beta+2\mathfrak{B}\beta\delta+\mathfrak{C}\delta\delta)&=a\mathfrak{C}'
\end{aligned}
\]
gilt. Wenn nicht $a=0$ ist, folgt offenbar, dass die Form $\mathfrak{F}$ in $\mathfrak{F}'$ durch die eigentliche Substitution $\alpha$, $\beta$, $\gamma$, $\delta$ übergeht. --- Nimmt man aber anstelle der ersten drei Gleichungen in $\Psi$ und $\Psi'$ die nächsten drei, so bestätigt man leicht drei den eben angegebenen völlig entsprechende Gleichungen, in denen anstelle des Faktors $a$ überall $b$ steht. Daraus sieht man, dass derselbe Schluss weiterhin gilt, sofern nicht $b=0$ ist. Nimmt man schließlich die letzten drei Gleichungen $\Psi$, $\Psi'$, so findet man ebenso, dass der Schluss gilt, wenn nicht $c=0$ ist. Da aber sicher nicht alle $a$, $b$, $c$ zugleich $=0$ sein können, geht notwendig die Form $\mathfrak{F}$ durch die Substitution $\alpha$, $\beta$, $\gamma$, $\delta$ in $\mathfrak{F}'$ über und ist ihr damit eigentlich äquivalent. \textit{Q. E. D.}

\Needspace{8\baselineskip}
\Article{241}
Eine Form wie $\mathfrak{F}$ oder $\mathfrak{F}'$, die entsteht, indem eine von drei gegebenen Formen mit der aus den beiden übrigen zusammengesetzten Form zusammengesetzt wird, nennen wir \textit{aus diesen drei Formen zusammengesetzt}. Nach dem vorigen Artikel kommt es dabei nicht darauf an, in welcher Reihenfolge die drei Formen zusammengesetzt werden. Ebenso: Sind beliebig viele Formen $f$, $f'$, $f''$, $f'''$ usw. gegeben (deren Determinanten im Verhältnis von Quadraten zueinander stehen müssen), und setzt man die Form $f$ mit $f'$, die resultierende mit $f''$, die daraus entstehende mit $f'''$ usw. zusammen, so heißt die am Ende dieser Operation entstehende Form \textit{aus allen Formen} $f$, $f'$, $f''$, $f'''$ usw. zusammengesetzt. Man beweist leicht, dass auch hier die Reihenfolge beliebig ist; \textit{d.h.}, in welcher Reihenfolge man die Formen auch zusammensetzt, die daraus entstehenden Formen sind stets eigentlich äquivalent. --- Weiter ist offenbar: Sind den Formen $f$, $f'$, $f''$ usw. jeweils die Formen $g$, $g'$, $g''$ usw. eigentlich äquivalent, so ist die aus diesen zusammengesetzte Form der aus jenen zusammengesetzten Form eigentlich äquivalent.
